\documentclass[acmsmall,nonacm]{acmart}
\pagestyle{empty}
\settopmatter{printacmref=false}
\begin{document}
\thispagestyle{empty}% acmart re-installs standardpagestyle \AtBeginDocument, overriding the preamble \pagestyle{empty}; with no \maketitle this would leak a page-1 folio. Re-assert empty so the titleless body-only fixture matches the Typst twin.

\section{Elements}
Text before a figure to establish context for the caption comparison below.

\begin{figure}[h]
\centering
\rule{4cm}{2cm}
\caption{A sample figure caption that is long enough to wrap onto a second line so we can examine caption typography.}
\label{fig:sample}
\end{figure}

\begin{table}[h]
\caption{A sample table caption set above the table body as ACM requires.}
\label{tab:sample}
\begin{tabular}{ll}
\toprule
Header A & Header B \\
\midrule
Value 1 & Value 2 \\
Value 3 & Value 4 \\
\bottomrule
\end{tabular}
\end{table}

\begin{theorem}\label{thm:main}
This is the statement of a theorem, which should be typeset with an italic body and a small-caps head.
\end{theorem}

\begin{lemma}
A lemma shares the theorem counter and uses the same plain style.
\end{lemma}

\begin{definition}
A definition uses an upright (roman) body and an italic head.
\end{definition}

\begin{proof}[Proof.]
This is a proof, ending with a QED square.
\end{proof}

\begin{itemize}
\item First bullet item in an itemized list.
\item Second bullet item, slightly longer so it may wrap around to a second line for spacing checks.
\end{itemize}

\begin{enumerate}
\item First enumerated item.
\item Second enumerated item.
\end{enumerate}

A lead paragraph precedes the display equation to anchor its vertical spacing.
\begin{equation}
a = b + c
\end{equation}
Text after the display equation continues at the margin with no added indent.

A lead paragraph precedes the verbatim block to anchor its vertical spacing.
\begin{verbatim}
code line one
code line two
\end{verbatim}
Text after the verbatim block continues at the margin with no added indent.

\appendix
\section{Supplementary material}
\begin{theorem}
A theorem inside the appendix should be numbered A.1, tracking the lettered section rather than an arabic one.
\end{theorem}

\end{document}
